\section{Introduction}
\label{sec:intro}

VLA policies map camera images, a language instruction and the robot state to actions with a pretrained vision-language model (VLM) and an action head~\cite{kim2024openvla,black2024pi0,shukor2025smolvla}.
Their cost is paid inside a control loop: at batch size one, every re-plan waits for the full model before the robot can act.
Work on efficient VLAs reduces this cost in several places: smaller backbones~\cite{shukor2025smolvla,wen2025tinyvla}, fewer tokens or layers per step~\cite{xu2025vlacache,yang2025efficientvla,yue2024deer}, cheaper action decoding~\cite{pertsch2025fast,kim2025oft}, asynchronous execution~\cite{black2025rtc,sendai2025a2c2}, and inference systems~\cite{ma2025realtime,islam2026onnx}.
Whether a given saving preserves closed-loop behaviour depends on the checkpoint and the hardware: exporting the SmolVLA checkpoint used here to ONNX cut its median latency by 40\% but, through truncated instructions, also cut LIBERO-Spatial success from 70\% to 41\%~\cite{islam2026onnx}.

This report examines SmolVLA~\cite{shukor2025smolvla}, a 0.45B-parameter flow-matching VLA designed for affordable hardware, on the LIBERO-Spatial benchmark~\cite{liu2023libero} and a Tesla T4, a GPU without bfloat16 tensor cores.
The scope is training-free: weights, architecture and benchmark protocol stay fixed, and only how the policy is executed changes.
The report covers the four parts of the assignment:
\begin{itemize}
\item a study of efficient VLAs and of diffusion versus flow-matching action generation, with a literature matrix of 18 verified sources (\cref{sec:study});
\item a reproduction of the released checkpoint with LeRobot, including the differences from the paper configuration (\cref{sec:reproduce});
\item an analysis of strengths, limitations, trajectories and computational bottlenecks (\cref{sec:analysis});
\item an improvement study that tests two hypotheses on identical episodes: re-planning less often, and re-planning more cheaply (\cref{sec:improve}).
\end{itemize}
The individual levers that make re-planning cheaper (half precision, fewer flow steps, CUDA graphs) are known~\cite{ma2025realtime,tang2026coda}; the contribution is a measured account of why each one helps on this hardware and a paired closed-loop test of whether the policy still succeeds.
